\chapter{Индуцированный фермионным детерминантом член Скирма}

Предыдущий гейт оставил один естественный источник общего действия:
проинтегрировать уже существующий фермионный пакет $H_{15}$ и получить
пространственные члены из
\begin{equation}
 \Gamma_{\rm F}[V]=-\log\det D[V].
 \label{eq:v5-fermion-determinant}
\end{equation}
Проверяется не формальная допустимость этого механизма, а более сильное
утверждение: выводит ли текущий родитель одновременно положительные
$E_2$, $E_4$ и их общий масштаб без нового массового параметра.

\section{Что действительно гарантирует литература}

Для четырёхмерного кирального оператора производное разложение нормальной
по чётности части детерминанта известно до четвёртого порядка в общей
матричной постановке; см. Л.~Л.~Сальседо,
\path{arXiv:hep-th/0012166}. Поэтому после задания оператора вида
\begin{equation}
 D_{M,V}=\mathord{\not\!\partial}
 +M\bigl(P_LV+P_RV^*\bigr)
 \label{eq:v5-chiral-dirac-coupling}
\end{equation}
можно получить структуру
\begin{equation}
 \Gamma_{\rm F}^{+}[V]
 =c_0+c_2\int\Tr(L_\mu L^\mu)
 +\sum_jc_{4,j}\int\mathcal O_{4,j}(L)+O(\partial^6),
 \qquad L_\mu=V^{-1}\partial_\mu V.
 \label{eq:v5-fermion-derivative-expansion}
\end{equation}
Один из инвариантов $\mathcal O_{4,j}$ является членом Скирма. Однако
общий результат содержит не только его. Уже классические вычисления
Аитчисона--Фрейзера и proper-time разложение в обзоре
\path{arXiv:hep-ph/9401216} дают рядом со стабилизирующим членом другой
четырёхпроизводный инвариант, который в части параметров дестабилизирует
малые профили. Обзор Д.~Дьяконова \path{arXiv:hep-ph/0009006} также
подчёркивает, что чистая модель Скирма является усечением: другие
инварианты четвёртого порядка имеют тот же производный порядок.

Следовательно, литература доказывает возможность индуцировать
четырёхпроизводный сектор, но не универсальную редукцию этого сектора к
единственному положительному $E_4$.

\section{Масса и регуляризация не сокращаются}

В proper-time схеме коэффициенты содержат неполные гамма-функции от
\begin{equation}
 x=\frac{M^2}{\Lambda^2}.
\end{equation}
Например, в явном четвёртом порядке встречаются одновременно
$\Gamma(1,x)$ и $\Gamma(2,x)$. Их отношение равно
\begin{equation}
 \frac{\Gamma(2,x)}{\Gamma(1,x)}=1+x,
 \label{eq:v5-incomplete-gamma-ratio}
\end{equation}
поэтому относительные веса четвёртых производных меняются при изменении
$M/\Lambda$. Квадратичный коэффициент имеет размер массы в квадрате,
тогда как коэффициенты четвёртого порядка безразмерны. Уже размерностно
радиус возможного баланса имеет вид
\begin{equation}
 R_*=\frac1M\,
 \mathcal R\!\left(\frac{M}{\Lambda},\{Y\},\text{схема}\right),
 \label{eq:v5-fermion-induced-radius}
\end{equation}
где $\{Y\}$ обозначает матрицы связи фермионов с фоном.

Работа Р.~Болла и Ж.~Рипки \path{arXiv:hep-ph/9312260} показывает, что
согласованная регуляризация вещественной и мнимой частей требует
дополнительного ультрафиолетового поведения фермионной самоэнергии.
Аномальная часть при этом может оставаться правильно нормированной, но
это не делает коэффициенты нормальной части независимыми от выбранного
массового профиля.

\begin{nogocriterion}[Топология не фиксирует нормальную часть]
Квантование аномальной фазы детерминанта не разрешает использовать тот же
аргумент для положительных коэффициентов $E_2$ и $E_4$. Фаза определяется
аномалией и спектральной асимметрией, а модуль детерминанта сохраняет
зависимость от массы, представления и перенормировки.
\end{nogocriterion}

\section{Сверка с пакетом \texorpdfstring{$H_{15}$}{H15}}

Проект уже доказал существование физических заряженных одноформ на
$H_{15}$, но не вывел единственный юкавский оператор: после фиксации
семейной грассмановой связности остаются две относительные амплитуды.
Отдельный гейт фермионной спектральной меры оставил независимыми чётный
момент, заряженно-лептонный коэффициент и масштаб $\Lambda$.

Поэтому формула \eqref{eq:v5-chiral-dirac-coupling} не является ещё одним
записыванием уже построенного $V_{15}$. Для неё недостают:
\begin{enumerate}
 \item каноническая карта от обменного KO7-унитария $V_{15}$ к
 фермионному массовому блоку на $H_{15}$;
 \item ненулевая фермионная щель $M$;
 \item фиксированные относительные юкавские амплитуды;
 \item условие перенормировки модуля детерминанта.
\end{enumerate}
Подстановка этих данных извне способна построить модель, но не является
выводом версии V.

Вес $1/7$ также не закрывает пробел. Если один и тот же нормированный след
умножает весь детерминант, он умножает и $E_2$, и все $E_4$ на общий
множитель. В отношении коэффициентов и в стационарном радиусе этот
множитель сокращается. Он сохраняет значение канального веса, но не
становится массой $M$.

\begin{theorem}[Условная индукция без замыкания масштаба]
Фермионный детерминант после внешнего задания $D_{M,V}$ действительно
может индуцировать квадратичный и четырёхпроизводный сектора. Но текущий
родитель не задаёт сам оператор связи, фермионную щель, относительные
амплитуды и схему модуля детерминанта. Кроме того, четвёртый порядок не
сводится автоматически к единственному положительному члену Скирма.
Следовательно, общий масштаб $E_2/E_4$ и конечный радиус из $H_{15}$,
$V_{15}$ и веса $1/7$ не выведены.
\end{theorem}

\section{Что остаётся содержательным}

Мнимая, нечётная по чётности часть детерминанта имеет более жёсткий
топологический статус. Работа \path{arXiv:hep-th/0012174} отделяет её от
нормальной части и связывает с аномальным функционалом. Поэтому следующий
допустимый тест должен быть не новой попыткой получить радиус, а проверкой
того, какую фазу Весса--Зумино--Виттена или эта-инвариант несёт обменная
Real-пара $V_{15}$ и сокращается ли она между ориентациями.

\begin{protocol}
\begin{enumerate}
 \item индуцирование $E_2$ фермионами: \textbf{условно допустимо};
 \item индуцирование инвариантов $E_4$: \textbf{условно допустимо};
 \item единственный положительный член Скирма: \textbf{не гарантирован};
 \item канонический оператор $D[V_{15}]$ на $H_{15}$:
 \textbf{не построен};
 \item фермионная щель $M$: \textbf{не выведена};
 \item независимость от $M/\Lambda$ и схемы: \textbf{нет};
 \item общий вес $1/7$ фиксирует радиус: \textbf{нет};
 \item физическое замыкание: \textbf{не достигнуто};
 \item следующий гейт: \textbf{аномальная фаза обменной Real-пары}.
\end{enumerate}
\end{protocol}

Машинный сертификат сохранён в
\path{s2t_v5_fermionic_determinant_induced_skyrme_gate_results.json}.